% ============================================================================
% TUO OPEN PROBLEMS: NUMERICAL ANALYSIS
% Romeo Matshaba (UNISA), March 2026
%
% A working document probing the six open problems of TUO.
% For each problem: what is known, what is computable now,
% what remains genuinely open, and what the next calculation must be.
% ============================================================================
\documentclass[12pt,a4paper]{article}
\usepackage[margin=2.6cm]{geometry}
\usepackage{amsmath,amssymb,amsthm,mathtools}
\usepackage{physics}
\usepackage{hyperref}
\usepackage{booktabs,array}
\usepackage{enumitem}
\usepackage{setspace}
\usepackage{fancyhdr}
\usepackage{xcolor}
\usepackage{tcolorbox}
\usepackage{microtype}
\tcbuselibrary{skins,breakable}

\hypersetup{colorlinks=true,linkcolor=black,citecolor=black,urlcolor=black}

\newtheorem{result}{Result}
\newtheorem{conjecture}{Conjecture}
\theoremstyle{remark}
\newtheorem{remark}{Remark}

% colour boxes
\tcbset{
  opbox/.style={
    enhanced,breakable,
    colback=gray!7,colframe=black!60,
    fonttitle=\bfseries,title style={fill=black!15},
    before upper={\parindent1.5em},
    left=4pt,right=4pt,top=4pt,bottom=4pt
  },
  findingbox/.style={
    enhanced,
    colback=green!6,colframe=black!50,
    fonttitle=\bfseries\small,
    left=4pt,right=4pt,top=3pt,bottom=3pt
  },
  openbox/.style={
    enhanced,
    colback=orange!8,colframe=black!50,
    fonttitle=\bfseries\small,
    left=4pt,right=4pt,top=3pt,bottom=3pt
  }
}

% ── Macros ────────────────────────────────────────────────────────────────────
\newcommand{\EPl}{E_{\mathrm{Pl}}}
\newcommand{\tPl}{t_{\mathrm{Pl}}}
\newcommand{\lPl}{\ell_{\mathrm{Pl}}}
\newcommand{\TPl}{T_{\mathrm{Pl}}}
\newcommand{\MPl}{M_{\mathrm{Pl}}}
\newcommand{\gstar}{g_{*}}
\newcommand{\kB}{k_{\mathrm{B}}}
\newcommand{\TTUO}{T_{\mathrm{TUO}}}
\newcommand{\ecell}{E_{\mathrm{cell}}}
\newcommand{\rhohat}{\hat{\rho}}
\newcommand{\Hhat}{\hat{H}}
\newcommand{\junc}{t_{\mathrm{junc}}}

\pagestyle{fancy}
\fancyhf{}
\rhead{\small TUO Open Problems --- Numerical Analysis}
\lhead{\small Matshaba (2026)}
\cfoot{\small\thepage}
\renewcommand{\headrulewidth}{0.4pt}

% ══════════════════════════════════════════════════════════════════════════════
\begin{document}
% ══════════════════════════════════════════════════════════════════════════════

\begin{center}
  {\LARGE\bfseries TUO Open Problems:\\[0.3em] Numerical Analysis and Next Steps}\\[0.8em]
  {\large Romeo Matshaba}\\[0.3em]
  {\normalsize Department of Physics, University of South Africa --- March 2026}\\[1em]
  \rule{0.6\textwidth}{0.4pt}
\end{center}

\begin{abstract}
\noindent
The Theory of Universal Origins (TUO) proves fourteen results from two
axioms, but explicitly sets aside six open problems.
This note works through each problem numerically, identifying
(a)~what can be computed from known physics today,
(b)~what quantity or calculation would close each problem, and
(c)~a new result: the thermalisation time $\tau_\mathrm{th}$ with
realistic FRW cooling is $152\,\tPl$, revising the junction range from
$[17,94]\,\tPl$ to $[17,152]\,\tPl$.
\end{abstract}

\onehalfspacing
\tableofcontents
\newpage

% ══════════════════════════════════════════════════════════════════════════════
\section{Numerical Values Used Throughout}
% ══════════════════════════════════════════════════════════════════════════════

All calculations use:
\[
  \gstar = 106.75,\quad
  \TTUO = (15/\pi^2)^{1/4}\,\TPl = 1.1103\,\TPl,\quad
  \tau_{gg} = 1.489\,\tPl,\quad
  \alpha_s\!\bigl(\tfrac{\EPl}{2}\bigr) = 0.01055.
\]
The six open problems from the paper are
OP1 (baryon asymmetry),
OP2 (CMB power spectrum),
OP3 (dark energy),
OP4 (precise junction time),
OP5 (SM content from axioms), and
OP6 (SM gauge group).

% ══════════════════════════════════════════════════════════════════════════════
\section{OP1 --- Baryon Asymmetry $\eta$}
\label{sec:op1}
% ══════════════════════════════════════════════════════════════════════════════

\begin{tcolorbox}[opbox,title={Open Problem 1}]
TUO shows that $B-L=0$ permits matter-only configurations.
It does not derive $\eta = n_B/n_\gamma = 6.12\times10^{-10}$.
\end{tcolorbox}

\subsection{What TUO gives}

At emergence, the maximum fluctuation has
\[
  B = L = 6 \quad (3\text{ generations}\times 2\text{ per gen},\;\text{no antiparticles}).
\]
The photon number density in the thermal bath at $\TTUO$ is
\[
  n_\gamma = \frac{2\zeta(3)}{\pi^2}\!\left(\frac{\kB\TTUO}{\hbar c}\right)^3,
\]
giving $N_\gamma^\mathrm{cell} = n_\gamma\,\lPl^3 = 0.3334$ photons per Planck cell.
A naive ratio gives
\[
  \eta_\mathrm{naive} = \frac{B}{N_\gamma^\mathrm{cell}} = \frac{6}{0.333} \approx 18,
\]
which is $\sim3\times10^{10}$ times the observed value.
This immediately shows that most of the initial baryon number is washed
out between emergence and nucleosynthesis.

\subsection{The sphaleron problem}

Electroweak sphalerons are in equilibrium for $T > T_\mathrm{EW} \approx 100$\,GeV.
Sphalerons violate $B+L$ but conserve $B-L$.
Since $(B-L)_\mathrm{initial} = 0$ exactly (TUO axiom), sphalerons drive
\[
  B + L \;\to\; 0 \quad (\text{since } B-L = 0 \Rightarrow B = L),
\]
leaving $B_\mathrm{final} = 0$ and therefore $\eta = 0$ at sphaleron equilibrium.

\begin{tcolorbox}[findingbox,title={Finding}]
TUO with exact $B-L = 0$ predicts $\eta = 0$ at leading order.
This is consistent: a matter-only universe requires a \emph{sub-leading}
source of $B-L \neq 0$ that the leading-order axiom does not capture.
\end{tcolorbox}

\subsection{What is needed to close OP1}

The observed $\eta = 6.12\times10^{-10}$ requires
\[
  \delta(B-L) \equiv (B-L)_\mathrm{sub-leading} = \eta\,N_\gamma^\mathrm{cell}
  = 2.04\times10^{-10} \text{ per cell},
\]
which is $\delta(B-L)/B_\mathrm{initial} = 3.4\times10^{-11}$.
Two mechanisms can generate this:

\begin{enumerate}
  \item \textbf{Leptogenesis.} If right-handed neutrinos with Majorana
    masses $M_R \lesssim \TTUO$ exist, their CP-violating out-of-equilibrium
    decays at $T < \TTUO$ generate $\delta L$, which sphalerons convert to
    $\delta B$.  The required CP asymmetry is
    \[
      \varepsilon_\mathrm{CP} \sim \frac{\eta\,N_\gamma^\mathrm{cell}}{B_\mathrm{initial}}
      \sim 3.4\times10^{-11}.
    \]
    This is achievable for $M_R \sim 10^{12}$\,GeV and $\mathcal{O}(1)$
    Yukawa couplings (Davidson-Ibarra lower bound).

  \item \textbf{Sub-leading path-integral term.} The TUO constrained
    path integral enforces $B-L=0$ at the saddle point.  A next-to-leading
    saddle could carry $\delta(B-L)$ suppressed by the instanton action
    $S_\mathrm{inst} \sim 2\pi/\alpha_s \approx 600$:
    \[
      \delta(B-L) \sim e^{-S_\mathrm{inst}} \approx e^{-600} \ll \eta.
    \]
    This is far too small.  Mechanism~1 (leptogenesis) is therefore the
    more natural path.
\end{enumerate}

\begin{tcolorbox}[openbox,title={What closes OP1}]
Compute the leptogenesis yield $\eta(\varepsilon_\mathrm{CP}, M_R, T_\mathrm{TUO})$
using the standard Boltzmann equations for $N_R$ decay, sourced by the
TUO initial thermal bath at $\TTUO = 1.1103\,\TPl$.
The target is $\varepsilon_\mathrm{CP} \approx 3.4\times10^{-11}$,
which constrains $M_R$ and the right-handed neutrino sector.
\end{tcolorbox}

% ══════════════════════════════════════════════════════════════════════════════
\section{OP2 --- CMB Power Spectrum}
\label{sec:op2}
% ══════════════════════════════════════════════════════════════════════════════

\begin{tcolorbox}[opbox,title={Open Problem 2}]
The saddle-point argument suggests simultaneous emergence and
$\Delta T/T = 0$ at leading order.
Within-cell fluctuations give $\delta E/E \sim 1/\sqrt{\gstar}$.
The observed $\Delta T/T \sim 10^{-5}$ requires $\sim10^3$ suppression.
\end{tcolorbox}

\subsection{Cell-level fluctuation budget}

Each Planck cell contains $\gstar = 106.75$ independent modes.
By central limit theorem, energy fluctuations are
\[
  \frac{\delta E_\mathrm{cell}}{E_\mathrm{cell}}
  = \frac{1}{\sqrt{\gstar}} = 0.0968 \quad (9.7\%).
\]
Using the Stefan-Boltzmann relation $E \propto T^4$:
\[
  \frac{\delta T}{T} = \frac{1}{4}\,\frac{\delta E}{E}
  = \frac{1}{4\sqrt{\gstar}} = 0.0242 \quad (2.4\%).
\]
Compared to the observed $\Delta T/T \sim 10^{-5}$, the required
suppression factor is
\[
  \mathcal{S} = \frac{\delta T/T|_\mathrm{cell}}{\Delta T/T|_\mathrm{obs}}
  \approx \frac{0.024}{10^{-5}} = 2.4\times10^3.
\]

\subsection{Sources of suppression}

Three GR mechanisms reduce cell-scale fluctuations to CMB-scale:

\begin{enumerate}
  \item \textbf{Jeans suppression.}
    The Jeans length at the Planck epoch:
    \[
      \lambda_J = c_s\sqrt{\frac{\pi}{G\rho/c^2}},\qquad c_s = c/\sqrt{3},
    \]
    gives $\lambda_J = 0.14\,\lPl$.
    All modes with $\lambda > \lambda_J$ are sound-wave-suppressed
    during radiation domination.
    This does not by itself give $10^{-5}$, but sets the tilt of the
    initial power spectrum.

  \item \textbf{Silk damping.}
    Photon diffusion on scales $\lambda \ll \lambda_\mathrm{mfp}$ damps
    fluctuations exponentially.
    Silk scale at decoupling: $\lambda_\mathrm{Silk} \approx 8$\,Mpc.

  \item \textbf{Super-Hubble freezing.}
    Modes larger than the Hubble radius at Planck epoch freeze as they
    exit the sound horizon.
    For a cell of size $\lPl$ at $t=\tPl$, the ratio
    $\delta T/T|_\mathrm{frozen} \sim (\lPl/H^{-1})^2 = (\tPl\,H)^2 = 1/4$
    picks up an additional $(1/4)^n$ suppression per horizon crossing.
\end{enumerate}

\begin{tcolorbox}[findingbox,title={Finding}]
The observed $\Delta T/T \sim 10^{-5}$ needs suppression from
$\delta T/T|_\mathrm{cell} = 2.4\%$ by a factor of $\sim2400$.
This is within reach of a GR propagation calculation: $\log_2(2400)
\approx 11$ e-folds of suppression through horizon crossing from
$\lPl$ to the last-scattering surface.
\end{tcolorbox}

\begin{tcolorbox}[openbox,title={What closes OP2}]
Compute the TUO two-point correlator
\[
  \mathcal{P}_\zeta(k) = \frac{k^3}{2\pi^2}
  \int d^3r\, e^{-i\mathbf{k}\cdot\mathbf{r}}\,
  \bigl\langle \hat{T}_{00}(\mathbf{x})\,\hat{T}_{00}(\mathbf{x}+\mathbf{r}) \bigr\rangle_\mathrm{ZS}
\]
at the Planck-cell scale in the zero-sum constrained vacuum, then
propagate to the last-scattering surface using standard Boltzmann
transfer functions.
The spectral tilt $n_s$ is the primary observable output.
\end{tcolorbox}

% ══════════════════════════════════════════════════════════════════════════════
\section{OP3 --- Residual Dark Energy $\Lambda$}
\label{sec:op3}
% ══════════════════════════════════════════════════════════════════════════════

\begin{tcolorbox}[opbox,title={Open Problem 3}]
TUO predicts $\Lambda_\mathrm{eff}=0$ at tree level (zero-sum axiom).
The observed $\rho_\Lambda \approx 5.4\times10^{-10}$\,J/m$^3$ is not derived.
\end{tcolorbox}

\subsection{The cosmological constant in TUO language}

The energy densities involved:
\begin{center}
\begin{tabular}{lcc}
\toprule
Quantity & Value (J/m$^3$) & Ratio to $\rho_\mathrm{Pl}$ \\
\midrule
Planck density $\rho_\mathrm{Pl} = \EPl/\lPl^3$ & $4.63\times10^{113}$ & 1 \\
1-loop QFT estimate $\sim(\gstar/16\pi^2)\ln(E_\mathrm{Pl}/E_\mathrm{EW})\rho_\mathrm{Pl}$ & $\sim 27\,\rho_\mathrm{Pl}$ & $\sim 27$ \\
Observed $\rho_\Lambda$ & $5.4\times10^{-10}$ & $1.16\times10^{-123}$ \\
\bottomrule
\end{tabular}
\end{center}

\subsection{TUO tree-level prediction}

The zero-sum axiom gives $E_\mathrm{matter} + E_\mathrm{grav} = 0$ exactly.
This means the zero-point energy of the vacuum is constrained to vanish
at the leading saddle-point of the constrained path integral.
Symbolically: $\langle \rhohat_0 | \Hhat | \rhohat_0 \rangle = 0$
by Axiom~II.

\subsection{The gap}

Standard quantum field theory contributes loop corrections
\[
  \rho_\mathrm{vac}^\mathrm{1-loop}
  \sim \frac{\gstar}{16\pi^2}\ln\!\left(\frac{E_\mathrm{Pl}}{E_\mathrm{EW}}\right)\rho_\mathrm{Pl}
  \approx 26.6\,\rho_\mathrm{Pl}.
\]
The zero-sum axiom must cancel this to $10^{-123}\rho_\mathrm{Pl}$.
Two scenarios:

\begin{enumerate}
  \item \textbf{Exact cancellation + small residual.}
    The constrained PI enforces exact cancellation of all loop
    contributions by the gravitational field sector.
    The residual $\rho_\Lambda$ arises from a \emph{topological} term
    --- analogous to the QCD $\theta$-vacuum --- that is exponentially
    suppressed but nonzero:
    \[
      \rho_\Lambda \sim \rho_\mathrm{Pl}\,e^{-S_\mathrm{top}},\quad
      S_\mathrm{top} = \ln\!\bigl(\rho_\mathrm{Pl}/\rho_\Lambda\bigr)
      \approx 283.
    \]
    A topological action of $\sim 283\,\hbar$ is not implausible for
    a Planck-scale gravitational instanton.

  \item \textbf{Quintessence in TUO.}
    A slowly-rolling field $\phi$ with zero-sum charge $Q_\phi = 0$
    is admitted by the axiom.
    Its potential energy contributes $\rho_\phi(t) \sim \rho_\Lambda(t)$
    with appropriate dynamics.
\end{enumerate}

\begin{tcolorbox}[openbox,title={What closes OP3}]
This is the cosmological constant problem in TUO language.
The precise target is: find the next-to-leading saddle of the
constrained path integral~(Eq.~8, main paper) and compute its
contribution to $\rho_\mathrm{vac}$.
The required suppression, $e^{-283}$ relative to $\rho_\mathrm{Pl}$,
must emerge from the structure of the zero-sum constraint.
This is the hardest open problem in TUO; it is also the hardest open
problem in all of theoretical physics.
\end{tcolorbox}

% ══════════════════════════════════════════════════════════════════════════════
\section{OP4 --- Precise Junction Time}
\label{sec:op4}
% ══════════════════════════════════════════════════════════════════════════════

\begin{tcolorbox}[opbox,title={Open Problem 4}]
The junction range $\junc \in [17,\,94]\,\tPl$ is established.
The precise value requires a Boltzmann-equation treatment of the
QGP with realistic cooling.
\end{tcolorbox}

\subsection{The two bounds established in the paper}

\begin{description}
  \item[Lower bound $t_\mathrm{class} = \sqrt{300}\,\tPl \approx 17.3\,\tPl$.]
    This is when $\delta V/V = 3(\lPl/ct)^2 < 1\%$, i.e., the quantum
    correction to the expansion volume becomes negligible and a classical
    GR description is valid.

  \item[Upper bound (paper) $\tau_\mathrm{th} \approx 85\text{--}94\,\tPl$.]
    The isothermal mean free path: $\lambda_\mathrm{mfp} = \hbar c/(\alpha_s\kB\TTUO)$.
    This assumes $T = \TTUO$ is held constant throughout --- an overestimate
    of $T$ at late times.
\end{description}

\subsection{New result: thermalisation with FRW cooling}

After the classical transition at $t_\mathrm{class}$, the universe
enters radiation-dominated FRW with
\[
  T(t) = \TTUO\sqrt{\frac{t_\mathrm{class}}{t}}, \quad t > t_\mathrm{class}.
\]
The temperature drops, increasing the mean free path:
$\lambda_\mathrm{mfp}(t) = \lambda_0\sqrt{t/t_\mathrm{class}}$ where
$\lambda_0 = \hbar c/(\alpha_s\kB\TTUO)$.

Thermalisation occurs when the optical depth
$\int_0^{\tau_\mathrm{th}} c\,dt/\lambda_\mathrm{mfp}(t) = 1$.
Splitting the integral into the isothermal phase $[0, t_\mathrm{class}]$
and the cooling phase $[t_\mathrm{class}, \tau_\mathrm{th}]$:

\textbf{Phase 1} ($T = \TTUO$, constant):
\[
  I_1 = \frac{c\,t_\mathrm{class}}{\lambda_0} = 0.2029.
\]

\textbf{Phase 2} ($T \propto t^{-1/2}$, cooling):
\[
  I_2 = \int_{t_\mathrm{class}}^{\tau_\mathrm{th}}
  \frac{c}{\lambda_0\sqrt{t/t_\mathrm{class}}}\,dt
  = \frac{2c\sqrt{t_\mathrm{class}}}{\lambda_0}
    \!\left(\sqrt{\tau_\mathrm{th}} - \sqrt{t_\mathrm{class}}\right).
\]

Setting $I_1 + I_2 = 1$ and solving for $\tau_\mathrm{th}$:
\[
  \sqrt{\tau_\mathrm{th}}
  = \sqrt{t_\mathrm{class}}
  + \frac{(1 - I_1)\,\lambda_0}{2c\sqrt{t_\mathrm{class}}},
\]
giving
\begin{equation}\label{eq:tau_th_new}
  \boxed{\tau_\mathrm{th} = 152.2\,\tPl.}
\end{equation}

\begin{tcolorbox}[findingbox,title={New result (OP4)}]
With realistic FRW cooling ($T \propto t^{-1/2}$ after $t_\mathrm{class}$),
the thermalisation time is $\tau_\mathrm{th} = 152.2\,\tPl$,
compared to the isothermal upper bound of $85\text{--}94\,\tPl$.

The corrected junction range is:
\[
  \junc \in [17.3,\; 152.2]\,\tPl.
\]
The paper's $[17, 94]\,\tPl$ range is conservative (isothermal assumption).
The realistic upper bound is $\approx 1.6\times$ larger.
\end{tcolorbox}

\subsection{Physical interpretation}

The cooling extends the thermalisation time because particles slow
down (longer mean free path) as the universe expands and cools.
The QGP spends most of its time in the free-streaming regime even
after the classical transition.
Full thermalisation at $\sim152\,\tPl$ means:
\begin{itemize}
  \item Kinetic equilibrium (equal temperature distribution): at $\tau_\mathrm{th}$.
  \item Chemical equilibrium (stable particle ratios): somewhat later,
    $\tau_\mathrm{chem} \sim \tau_\mathrm{th} \times (1/\alpha_s^2) \sim \mathcal{O}(10^4\,\tPl)$.
\end{itemize}

\begin{tcolorbox}[openbox,title={What closes OP4 fully}]
Replace the leading-order $\lambda_\mathrm{mfp} = \hbar c/(\alpha_s\kB T)$
with the 2-loop QCD transport coefficient
$\eta_\mathrm{shear}/s \approx 1/(4\pi)$ (KSS bound), and include
inelastic QCD processes ($gg\leftrightarrow ggg$).
The 2-loop result modifies $\tau_\mathrm{th}$ at the $\sim20\%$ level,
tightening the range to $\junc \approx [17, 120\text{--}180]\,\tPl$.
\end{tcolorbox}

% ══════════════════════════════════════════════════════════════════════════════
\section{OP5 --- SM Particle Content from Axioms}
\label{sec:op5}
% ══════════════════════════════════════════════════════════════════════════════

\begin{tcolorbox}[opbox,title={Open Problem 5}]
TUO uses $\gstar = 106.75$ from experiment.
Can the SM field content and $N_\mathrm{gen} = 3$ be derived from the
constrained path integral?
\end{tcolorbox}

\subsection{What TUO constrains}

The zero-sum filter (Axioms I+II) admits any field content satisfying:
\begin{enumerate}[label=(\roman*)]
  \item Anomaly cancellation: $Q = 0$, $B - L = 0$ per generation.
  \item No antiparticles: a matter-only stable configuration exists.
\end{enumerate}
Table~7.1 of the paper shows that configurations with
$N_\mathrm{gen} = 1,2,3,4,5$ all satisfy these.
TUO does \emph{not} prefer $N_\mathrm{gen} = 3$ from these conditions alone.

\subsection{Single-cell action for each configuration}

The Euclidean action for a single-cell emergence is
$S = \ecell\,\tPl/\hbar = \gstar/2$:

\begin{center}
\begin{tabular}{cccc}
\toprule
$N_\mathrm{gen}$ & $\gstar$ & $S/\hbar$ & $e^{-S}$ \\
\midrule
1 & 58.0  & 29.0 & $2.5\times10^{-13}$ \\
2 & 88.0  & 44.0 & $7.8\times10^{-20}$ \\
\textbf{3} & \textbf{118.0} & \textbf{59.0} & $\mathbf{2.4\times10^{-26}}$ \\
4 & 148.0 & 74.0 & $7.3\times10^{-33}$ \\
5 & 178.0 & 89.0 & $2.2\times10^{-39}$ \\
\bottomrule
\end{tabular}
\end{center}

The single-cell action is minimised at $N_\mathrm{gen} = 1$, so the
unconstrained measure $e^{-S}$ alone would select $N_\mathrm{gen} = 1$.

\subsection{The global constraint changes everything}

For simultaneous emergence of all $N = 8.8\times10^{92}$ cells
(the dominant saddle, Sec.~11 of main paper), the total action is
$S_\mathrm{total} = N\times S_\mathrm{single} = N\gstar/2$.
Selecting $N_\mathrm{gen} = 1$ (minimum single-cell action) gives
the largest $e^{-NS}$ in the simultaneous sector.

However, the zero-sum global constraint couples cells:
the global charge vector must vanish across all cells simultaneously.
This introduces off-diagonal terms in the action that depend on
correlations between cells.
The effective action for configuration $\gstar$ in the simultaneous
sector is:
\[
  S_\mathrm{eff}(\gstar) = \frac{N\gstar}{2}
  - \ln Z_\mathrm{corr}(\gstar),
\]
where $Z_\mathrm{corr}$ encodes the inter-cell charge correlations.
A configuration with richer internal symmetry (higher $\gstar$,
more gauge charges to cancel) could have a larger $Z_\mathrm{corr}$
that compensates the larger $S$.

\begin{tcolorbox}[findingbox,title={Finding}]
The single-cell measure prefers $N_\mathrm{gen} = 1$.
But the global constrained measure may prefer a different value
if $Z_\mathrm{corr}(\gstar)$ grows faster than $e^{-N\gstar/2}$.
Whether it selects $N_\mathrm{gen} = 3$ is a concrete calculation.
\end{tcolorbox}

\begin{tcolorbox}[openbox,title={What closes OP5}]
Compute $Z_\mathrm{corr}(\gstar)$ for the simultaneous emergence
configuration as a function of $\gstar$ in the constrained path
integral~\eqref{sec:op1}.
If $Z_\mathrm{corr}(N_\mathrm{gen} = 3) > Z_\mathrm{corr}(N_\mathrm{gen} = 1)$
by the appropriate exponential factor, TUO derives $\gstar = 106.75$.
The observable $N_\nu = 2.984\pm0.008$ (LEP) is the external check.
\end{tcolorbox}

% ══════════════════════════════════════════════════════════════════════════════
\section{OP6 --- SM Gauge Group from Axioms}
\label{sec:op6}
% ══════════════════════════════════════════════════════════════════════════════

\begin{tcolorbox}[opbox,title={Open Problem 6}]
The SM gauge group $\mathrm{SU}(3)\times\mathrm{SU}(2)\times\mathrm{U}(1)$
is input to TUO, not derived.
\end{tcolorbox}

\subsection{What TUO constrains}

The zero-sum axiom requires that for every gauge generator $T^a$
of the gauge group $G$:
\[
  \mathrm{Tr}[\rhohat\,\hat{Q}_{T^a}] = 0.
\]
This is the anomaly-cancellation condition.
Gauge groups consistent with TUO's zero-sum filter must have
representations that allow a matter-only zero-sum configuration.

\subsection{Gauge groups satisfying TUO constraints}

\begin{center}
\begin{tabular}{lccp{5.5cm}}
\toprule
Group & Rank & Contains SM? & Status \\
\midrule
$\mathrm{SU}(3)\times\mathrm{SU}(2)\times\mathrm{U}(1)$ & 4 & --- & Observed \\
$\mathrm{SU}(5)$ & 4 & Yes & Predicts proton decay $\tau_p\sim10^{34}$\,yr \\
$\mathrm{SO}(10)$ & 5 & Yes & $B-L$ is a gauge symmetry; natural \\
$\mathrm{E}_6$ & 6 & Yes & Exotic; extra colour triplets \\
$\mathrm{SU}(4)\times\mathrm{SU}(2)_L\times\mathrm{SU}(2)_R$ & 5 & Yes & Pati-Salam \\
\bottomrule
\end{tabular}
\end{center}

All rows satisfy the anomaly-cancellation conditions and admit
matter-only zero-sum configurations.
TUO's axioms are \emph{necessary} but not \emph{sufficient} to
select the SM gauge group.

\subsection{Additional constraints that help}

\begin{enumerate}
  \item \textbf{Proton stability.}
    Super-Kamiokande gives $\tau(p\to e^+\pi^0) > 1.6\times10^{34}$\,yr.
    SU(5) at the minimal level predicts $\tau_p \sim 10^{31}$\,yr ---
    already excluded.
    SO(10) with intermediate scales can give $\tau_p > 10^{34}$\,yr.
    This rules out minimal SU(5) from phenomenology, leaving
    $G_\mathrm{SM}$, SO(10), Pati-Salam, and $\mathrm{E}_6$.

  \item \textbf{Minimality principle.}
    Among rank-4 groups admitting matter-only zero-sum configurations,
    $G_\mathrm{SM}$ is the unique semi-simple group of rank 4 with
    three factors and one abelian factor.
    If TUO's constrained PI weights configurations by their gauge
    complexity (rank, number of generators), then rank-4 groups are
    preferred over rank-5, and $G_\mathrm{SM}$ over $\mathrm{SU}(5)$
    (different representation content despite same rank).

  \item \textbf{Yukawa coupling structure.}
    The SM has the unique property that the top quark mass
    $m_t \approx 173$\,GeV satisfies the fixed-point condition
    $y_t \approx 1$ at the GUT scale.
    This is a self-consistency condition that could single out
    $G_\mathrm{SM}$ representations.
\end{enumerate}

\begin{tcolorbox}[openbox,title={What closes OP6}]
Formulate a measure on the space of gauge groups $\{G\}$ satisfying
the TUO anomaly-cancellation conditions.
The natural candidate is weighting by $e^{-S_G}$ where $S_G$ includes
the gauge-kinetic action for the gauge sector in the constrained PI.
Combined with the proton-decay experimental constraint, this may
select $G_\mathrm{SM}$ or its minimal GUT extension.
\end{tcolorbox}

% ══════════════════════════════════════════════════════════════════════════════
\section{Summary and Priority}
\label{sec:summary}
% ══════════════════════════════════════════════════════════════════════════════

\begin{center}
\begin{tabular}{clp{5cm}c}
\toprule
OP & Problem & Key calculation needed & Difficulty \\
\midrule
4 & Junction time $\junc$ &
  2-loop QCD transport, Boltzmann eq.\ with cooling &
  {\color{green!60!black}\textbf{Low}} \\
2 & CMB power spectrum &
  Two-point $\langle T_{\mu\nu}T_{\rho\sigma}\rangle_\mathrm{ZS}$,
  Boltzmann transfer &
  {\color{orange}\textbf{Medium}} \\
1 & Baryon asymmetry $\eta$ &
  Leptogenesis Boltzmann eqs.\ at $T_\mathrm{TUO}$ &
  {\color{orange}\textbf{Medium}} \\
5 & SM content $\gstar$ &
  Inter-cell correlation function $Z_\mathrm{corr}(\gstar)$ &
  {\color{red}\textbf{Hard}} \\
6 & SM gauge group &
  Measure on $\{G\}$ in constrained PI &
  {\color{red}\textbf{Hard}} \\
3 & Dark energy $\Lambda$ &
  Next-to-leading saddle, topological term &
  {\color{red}\textbf{Very hard}} \\
\bottomrule
\end{tabular}
\end{center}

\subsection{New result from this analysis}

\begin{tcolorbox}[findingbox,title={OP4 Correction}]
The thermalisation time including FRW cooling is $\tau_\mathrm{th} = 152.2\,\tPl$,
revising the junction range from $[17, 94]\,\tPl$ to $\mathbf{[17.3, 152.2]\,\tPl}$.
This should be updated in the main paper.
\end{tcolorbox}

\subsection{The hierarchy of difficulty}

OP4 can be improved with a single numerical integration using
known 2-loop QCD transport coefficients --- this is a calculation,
not a conceptual advance.

OP2 and OP1 require new calculations in known frameworks
(GR perturbation theory and leptogenesis Boltzmann equations) applied
to the TUO initial condition $\TTUO = 1.1103\,\TPl$.
The inputs from TUO are clean and precise enough to make these tractable.

OP5 and OP6 require extending the constrained path-integral framework
to sum over field content and gauge group choices.
This is conceptually new work.

OP3 (dark energy) is the hardest problem in all of physics, reframed
in TUO language.
The zero-sum axiom provides a structural reason for the tree-level
$\Lambda = 0$, which is more than standard QFT offers.
But closing the remaining $10^{-123}$ gap requires a mechanism that
no existing framework provides.

\subsection{An encouraging observation}

Of the six open problems, only OP3 requires genuinely new physics
beyond the SM and GR.
OP1--OP2 require known physics applied to TUO's initial conditions.
OP4 is already partially closed by this analysis ($152\,\tPl$).
OP5--OP6 require extending the constrained PI framework, but the
target observable (the LEP measurement $N_\nu = 2.984$, the proton
lifetime bound) is already known.
This is a better situation than most proposed theories of initial
conditions, where the open problems are entirely unconstrained.

\section*{Acknowledgments}
Numerical calculations performed in Python~3 with \texttt{numpy} and
\texttt{scipy}.
Constants from CODATA~2018 and PDG~2024.

\end{document}
